\scanpage{5}
\[
\operatorname{réal}(S,A,R)\lhook\joinrel\longrightarrow E
\]
telle que l'image d'un arc soit un segment.

\uline{Prop} Pour qu'un graphe fini admette une réalisation géom. rectiligne, (dans un vectoriel convenable sur $\mathbb{R}$) il f et s qu'il \deleted{\unsure{a|est}} soit strict -- et alors il admet une réalisation rectiligne dans $\mathbb{R}^{3}$. De façon précise, si $S\hookrightarrow\deleted{\unsure{\mathbb{R}^{3}|R^{3}}}E$ \inserted{$(\simeq\mathbb{R}^{3},\ \unsure{v.\,ex.|p.\,ex.|v.a.})$} est une immersion telle que
% The pencil drawing stands in the left margin beside this condition.
\margin{\figure{\begin{tikzpicture}[x=.012cm,y=.012cm,line cap=round,line join=round]
\draw[line width=.65pt] (3,4)--(77,84);
\draw[line width=.6pt,rotate around={-25:(0,0)}] (0,0) ellipse[x radius=4.1,y radius=2.8];
\draw[line width=.55pt] (-.7,-3)--(-1.2,-6);
\fill (76,83)--(80,84.5)--(80.5,89)--(77,91)--(76.2,88)--(74.5,86.5)--cycle;
\draw[line width=.2pt] (87,95)--(148,158);
\fill (145,157)--(149.4,157)--(149,161.5)--(146.8,161)--cycle;
\draw[line width=.2pt] (147.5,159)--(149,163);
\end{tikzpicture}}}
\[
\forall x,y\in S\quad x\ne y,\quad \text{on ait}\quad
\deleted{\unsure{\{|\exists}}
\underbrace{\overline{xy}}_{\substack{\text{segm}\\\text{joignant}\\x\ \text{à}\ y}}\cap S=\{x,y\},
\]
\inserted{(i.e. trois points \inserted{distincts} jamais alignés)}
\inserted{et $\forall x,y,z,t$ distincts, $(\overline{x,y})\cap(\overline{z,t})=\varnothing$}

$\exists 1$ réalis géom. rectiligne $\unsure{\operatorname{réll}|\operatorname{réel}|\operatorname{réal}}(S,A,R)\lhook\joinrel-E$ qui $f$ \unsure{prolonge\textsuperscript{?}|prolonge}\ldots
% The possible author's question mark is over the l of prolonge.

Graphe connexe $=$ réal. géom \unsure{connex|connexe} $=$ qui n'est pas \deleted{\unsure{réunion}} ``somme'' de deux autres graphes $\ne\varnothing$

\uline{NB} Tout graphe est de façon unique réunion de graphes connexes \inserted{et disjoints}

[ graphe somme : $S\amalg S',\quad A\amalg A',\quad R\amalg R'\ \ldots$ ]

% A curved arrow in the left margin runs from beside NB to Graphe.
\margin{\figure{\begin{tikzpicture}[baseline=-.5ex,x=1cm,y=1cm]
\draw[->] (.85,1.55) .. controls (.65,1.68) and (-.08,.62) .. (.12,.25) .. controls (.3,.02) and (1.02,.14) .. (1.55,0);
\end{tikzpicture}}}
% multi is encircled above polygonal. The parenthesis is added above the line;
% the bracketed condition shares the baseline of Graphe ... polygonal.
\[
\begin{aligned}
&\qquad\qquad\text{\inserted{(ou \uline{contour} (combinatoire))}}\\[-.3ex]
&\text{Graphe \inserted{multi}polygonal}\quad
\bigl[\quad\forall s\in S,\quad\operatorname{card}R\{s\}=2\quad\bigr]
\end{aligned}
\]

Dualité des \deleted{graphes polygonaux} contours (combinatoires).

\deleted{\unsure{donq|donc}} les composantes \inserted{\unsure{conn|conn.|connexes}} d'un \deleted{graphe polygonal}\inserted{contour} sont \deleted{polygonales.} des contours.

Classification des \deleted{graphes polygonaux}\inserted{contours} conn. $\ne\varnothing$

\uline{Prop} Un \deleted{graphe polyg.}\inserted{contour} connexe \inserted{$\ne\varnothing$} est connu à isom. près quand on connait $\operatorname{card}S$ ($=\operatorname{card}A$), qui est n'importe quel cardinal \deleted{fini ou infini}
\[
2\le c\le\aleph_{0}
\]
(strict \unsure{sss|ssi} $c\ne2$ i.e. $c\geqslant3$)
% archive sheet number: 4
